\scanpage{4}
% archive foliation: 3 (upper right and lower left)
Alors on a, \deleted{si $X$ purement de dim $n$ et $i\leq n$ (resp.\ $i\geq n$)}
\[
\begin{aligned}
\operatorname{Filt}^{j}H^i(\overline X,\mathbb{Q}_\ell)
={}&\text{partie de }T^\ell(h^i(X))\ \deleted{\text{provenant}}\inserted{\text{image}}\text{ du}\\
&\text{plus grand sous-motif de}\\
&\text{niveau }\leq i-2j\quad
\deleted{(\text{resp. }\unsure{2n-i-2j|2n-i-j})}\ 
\deleted{\text{de}}\inserted{\text{de}}\ h^i(X)
\end{aligned}
\]

Cette caractérisation \inserted{du niveau des motifs} est \unsure{tout} au moins correcte si on admet la résolution des singularités, \deleted{sur $k$}\inserted{des variétés de dim $\leq\dim X=n$ sur $k$.}
Comme conséquence de la résolution et des conjectures \unsure{standards|standarts}, on trouve donc que les dimensions des $\operatorname{Filt}^jH^i(\overline X,\mathbb{Q}_\ell)$ sont indépendantes de $\ell$.
(Bien mieux, ils correspondent à des motifs, et les \unsure{petits} facteurs successifs sont \deleted{\unsure{et}} purs de poids $i-2j$ resp $\unsure{2n-i-j|2n-i-2j}$).
\margin{Parler des inv. birationnels\dots}
\margin{Conjectures de Tate et de Hodge pourraient \uuline{passer ici}.}

7. \uline{Structure des $\otimes$-catégories}. \deleted{\unsure{abéliennes}}

\deleted{\unsure{Catégorie}} Catégorie abélienne avec opération $\otimes$ ``ass.'', ``comm.'' et ``unitaire''. On suppose de plus l'existence de \uline{Hom} internes, au sens de $\otimes$.
On dit qu'elle est ``rigide'' si
\[
\check M\otimes N\xrightarrow{\sim}\underline{\operatorname{Hom}}(M,N)
\]
Supposons de plus que $\mathcal C$ soit $K$-linéaire ($K$ un corps)
(\unsure{revient} à donner $K\longrightarrow\operatorname{End}(\operatorname{id}_{\mathcal C})$), et que les Hom soient de dim finie sur $k$ : \deleted{\unsure{ce est}} On dira qu'on a une $\otimes$-catégorie sur $k$ (pour simplifier).
Une telle catégorie est dite ``connexe'' si
\[
\operatorname{End}(\mathbf1_{\mathcal C})\simeq K\qquad(\unsure{par}\text{ l'identité}).
\]
